% !TEX program = xelatex
\documentclass[titlepage=formal,code=listings]{../SUSTechHomework}

\title{Project Report Example}
\coursecode{EE202}
\coursename{Embedded Systems Design}
\semester{Spring 2026}
\instructor{Prof. Example}
\reporttype{Course Project Report}
\date{\today}
\addbibresource{examples/project-report.bib}

\begin{document}

\clearauthors
\addauthor{Alice Example}{12340001}{alice@example.com}{Department of Electronic and Electrical Engineering}{Zhixin College}{System Design}
\addauthor{Bob Example}{12340002}{bob@example.com}{Department of Computer Science and Engineering}{Shuren College}{Software Integration}

\clearcontributions
\addcontribution{Alice Example}{Conceptualization, Hardware, Investigation}
\addcontribution{Bob Example}{Software, Validation, Writing - Original Draft}

\clearcredits
\addcredit{Hardware}{Alice Example implemented the sensing and acquisition pipeline.}
\addcredit{Software}{Bob Example integrated the control logic and visualization workflow.}

\maketitle
\tableofcontents

\section{Introduction}

This report demonstrates a complete project-style document with figures, code,
math, references, and structured team metadata.

\section{System Model}

The controller state is modeled as
\begin{align}
  \vect{x}_{k+1} &= \mat{A}\vect{x}_k + \mat{B}\vect{u}_k, \\
  \vect{y}_k &= \mat{C}\vect{x}_k.
\end{align}

The sampling interval is \qty{5}{\milli\second}, and the nominal current limit
is \qty{1.5}{\ampere}.

\section{Implementation}

\begin{sustechcodeblock}{C++}[Controller Loop][src/controller.cpp][Main project control loop.][code:controller]
for (int step = 0; step < max_steps; ++step) {
    update_state();
    apply_control();
}
\end{sustechcodeblock}

The control loop in \sustechcoderef{code:controller} follows the structure
recommended by \textcite{franklin-feedback}.

\section{Results}

\begin{sustechfiguregroup}
  \begin{sustechsubfigure}[0.48\linewidth][Baseline]
    \sustechplaceholdergraphic[0.92\linewidth][3cm]
  \end{sustechsubfigure}
  \sustechfloatsep
  \begin{sustechsubfigure}[0.48\linewidth][Improved]
    \sustechplaceholdergraphic[0.92\linewidth][3cm]
  \end{sustechsubfigure}
  \caption{Controller performance comparison}
  \label{fig:project-results}
\end{sustechfiguregroup}

As shown in \cref{fig:project-results}, the tuned controller reduces overshoot.
The reporting workflow also benefits from package guidance such as
\parencite{ctan-biblatex}.

\section{Conclusion}

This project example is suitable as a starting point for lab reports, course
projects, and team-based engineering writeups.

\sustechprintbibliography

\end{document}
